(1) 在忽略导体表面粗糙度、外部泄漏和高阶模的TEM近似下，圆柱对称高斯面给 $E_r(r)=q'/(2\pi\varepsilon r)$，$a<r<b$。环路安培定律给 $H_\phi(r)=I/(2\pi r)$，相应 $B_\phi=\mu I/(2\pi r)$。内外导体分别载反向电流，理想外导体外场相互抵消。
(2) 两导体的电势差 $V=\int_a^bE_rdr=q'\ln(b/a)/(2\pi\varepsilon)$，故 $C'=q'/V=2\pi\varepsilon/\ln(b/a)$。电能每单位长度为 $\frac12\varepsilon\int_a^b E_r^2(2\pi r\,dr)=q'^2/(2C')$。同样磁能每单位长度为 $\frac12\mu\int_a^bH_\phi^2(2\pi r\,dr)=L'I^2/2$，得到 $L'=\mu\ln(b/a)/(2\pi)$。
(3) 无损电报方程为 $\partial_xV=-L'\partial_tI$、$\partial_xI=-C'\partial_tV$，联立得行波速度 $v=1/\sqrt{L'C'}=1/\sqrt{\mu\varepsilon}$。特征阻抗 $Z_0=\sqrt{L'/C'}=[\ln(b/a)/(2\pi)]\sqrt{\mu/\varepsilon}$，单位为欧姆且随介质磁电参数改变；乘积 $L'C'=\mu\varepsilon$ 与几何对数相互抵消。
(4) 单向传播波满足 $V=Z_0I$，故每单位长度电场能 $u_E=C'V^2/2$、磁场能 $u_B=L'I^2/2$。因为 $C'Z_0^2=L'$，两者相等，$u_E=u_B$；波的能量传输率 $v(u_E+u_B)=vL'I^2=Z_0I^2=VI$，与沿横截面积分 $\boldsymbol E\times\boldsymbol H$ 的坡印廷能流一致。
